\documentclass[svgnames]{article}     % use "amsart" instead of "article" for AMSLaTeX format
%\geometry{landscape}                 % Activate for rotated page geometry

%\usepackage[parfill]{parskip}        % Activate to begin paragraphs with an empty line rather than an indent

\usepackage{graphicx}                 % Use pdf, png, jpg, or eps§ with pdflatex; use eps in DVI mode

%maths                                % TeX will automatically convert eps --> pdf in pdflatex
\usepackage{amssymb}
\usepackage{amsmath}
\usepackage{esint}
\usepackage{geometry}
\usepackage{fontawesome5}
\usepackage{hyperref}
\hypersetup{
    colorlinks=true,
    linkcolor=magenta,
    filecolor=magenta,
    urlcolor=magenta,
    }

% Inverting Color of PDF
%\usepackage{xcolor}
%\pagecolor[rgb]{0.19,0.19,0.19}
%\color[rgb]{0.77,0.77,0.77}

%noindent
%\setlength\parindent{0pt}

%pgfplots
\usepackage{pgfplots}

%images
\graphicspath{{/Users/devaldeliwala/research/paper_presentation/ocf_writeup/images/}}                   % Activate to set a image directory

%tikz
\usepackage{pgfplots}
\pgfplotsset{compat=1.15}
\usepackage{comment}
\usetikzlibrary{arrows}
\usepackage[most]{tcolorbox}

%Figures
\usepackage{float}
\usepackage{caption}
\usepackage{lipsum}


\title{Optimized Curve Fitting (OCF) Algorithm}
\author{Deval Deliwala}
\date{}                              % Activate to display a given date or no date

\begin{document}
\maketitle
%\section{}
%\subsection{}
%\tableofcontents                     % Activate to display a table of contents

\begin{center}
  GOAL \\ \vspace{5px}
\end{center}

The Optimized Curve Fitting (OCF) algorithm aims to determine the slope of the
Red Clump (RC) cluster in a color-magnitude diagram (CMD) with minimal manual
intervention. Traditional methods often rely on visual adjustments and unsharp
masking techniques to identify the RC slope, leading to potential
inconsistencies and biases. The OCF algorithm provides a systematic and
reproducible approach to identifying the RC slope by segmenting the RC bar into
defined tiles and iteratively fitting compound models on the extracted star data.

\section{Procedure}

\subsection{Inputs \& Parameters}

The algorithm requires the following inputs: 

\begin{itemize}
  \item[1.] \textbf{Star Catalogs}:\\ Two matched star catalogs, designated as
    \verb|catalog1| and \verb|catalog2|, are used to generate the CMD. The CMD
  is plotted as \verb|catalog1| $-$ \verb|catalog2| vs. \verb|catalogy|,
  where \verb|catalogy| is either \verb|catalog1| or \verb|catalog2|. 
\item[2.] \textbf{Initial Parallel Cutoffs}:\\ A set of parallel lines that
  \textit{roughly} delineate the RC slope. These lines do not need to be highly
  precise but should tightly bound the RC bar. 
\item[3.] \textbf{$x$-Range Cutoff}:\\ The range in the $x$-axis color space that
  defines the width of the RC bar. 
\item[4.] \textbf{Tile Count $n$}:\\ An integer specifying the number of
  segments/tiles to divide the RC bar. 
\end{itemize}

The provided inputs for an example CMD are visually represented in Figure 1, illustrating the
parallel and $x$-range cutoff used to bound the RC cluster. 

\begin{figure}[H]
  \centering
    \includegraphics[width = \textwidth]{F115W-F212N-vsF115W_cutoff.png}
    \caption{Example NRCB1 $F115W - F212N$ vs. $F115W$ provided cutoffs}
\end{figure}

\subsection{Preliminary Steps}

\subsubsection{Extend Parallel Cutoffs} 

The algorithm begins by extending the width of the initial parallel lines by
a factor of 3, allowing for a broader range of stars to be included in the
analysis. The expanded lines are also shown in Figure 1 in blue.

\subsubsection{Generate the First Tile} 

Using the extended parallel lines and the specified tile count $n$, the
algorithm generates a \textit{starting tile} of width \verb|x_range| / $n$. It
is placed at the left-most end of the provided \verb|x-range|, in between the parallel cutoffs,
extending horizontally by the calculated tile width in a left-riemann
fashion. An example starting bin for $n = 9$ is shown in Figure 2 below. 

\begin{figure}[H]
  \centering
    \includegraphics[width = 9cm]{starting_bin.png}
    \caption{Starting bin for $n = 9$} 
\end{figure}

\subsubsection{Extract Stars from Starting Tile} 

The algorithm then extracts the indices of stars that fall within the bounds of
the first tile. The corresponding magnitudes from \verb|catalog1| and
\verb|catalog2| are retrieved and these stars are compiled into a pandas
DataFrame object for subsequent analysis. 

\subsection{OCF Algorithm} 

Once the magnitudes of the stars in the starting bin have been extracted, teh
algorithm proceeds to the OCF phase for this specific bin. The primary goal of
OCF is to produce a compound model using the \verb|astropy.models| library that
fits the hitogram of \verb|catalogy| magnitude data optimally, thereby
minimizing the error of the mean magnitude for stars in that bin. 

The compound model used consists of a Gaussian component (\verb|Gaussian1D()|)
and a linear component (\verb|Linear1D()|). This combination is chosen because
it allows the mean magnitude to be calculated more accurately by focusing the
Gaussian fit on the dense region of stars, while the linear component accounts
for outliers and stars outside the dense region. 

\subsubsection{Compound Model Parameters} 

A compound \verb|Gaussian1D() + Linear1D()| model requires five parameters: 

\begin{itemize}
  \item[1.] Gaussian Component: 
    \begin{itemize}
      \item[-] Amplitude
        \item[-] Mean 
          \item[-] Standard Deviation
    \end{itemize}
  \item[2.] Linear Component: 
    \begin{itemize}
      \item[-] Slope
        \item[-] Intercept
    \end{itemize}
\end{itemize}

The objective of OCF is to determine the optimal values for these parameters
for a compound model such that the error in the model's returned mean parameter
is minimized.

\subsubsection{OCF Steps}

The OCF Algorithm employs two main iterative steps to achieve optimal fitting: 

\begin{itemize}
  \item[1.] \textbf{Looser-Requirement Iterative Compound Fitting}:\\
  - This step aims to determine suitable initial parameter values for the
  Gaussian mean and standard deviation. It uses the overall mean and standard
  deviation from the generic \\ \verb|scipy.stats.norm.fit()| fit as starting
  parameters for the compound models. 

  - The algorithm iterates through trial compound fits of increasing amplitude,
  accepting every fit that meet loose requirements for the returned standard
  deviation and amplitude. These fits help to calculate a more precise mean and
  standard deviation, which are used as the input parameters for the compound
  models in the next step. 

  \item[2.] \textbf{Strict Requirement Iterative Compound Fitting}:\\ 
  - Using the mean and standard deviation parameters from the previous step as
  more precise model input parameters, this phase generates new compound models
  and again, iterates through various increasing amplitude values. 

  - The algorithm employs much more strict requirements on the returned fitting
  values, ensuring the returned parameters (mean, standard deviation,
  amplitude) fall within acceptable ranges that ensure the fit was a success.

  - In addition to the returned fitting parameter requirements, a synthetic
  dataset is also generated to produce a histogram that similarly follows the
  shape of compound model. The stricter requirements also require
  a \verb|ks_statistic| and an earth mover's distance \verb|EMD_metric| to fall
  within acceptable ranges that guarantee the model's goodness-of-fit. 

  - Parameters from successful fits are stored for further analysis. 
\end{itemize}

\subsubsection{Error Calculation} 

If a compound model passes the strict requirements, the error of the mean is
calculated using one of two methods, allowing the user to choose. 

\begin{itemize}
  \item[1.] \textbf{Bootstrapping}: \\
  This method is computationally intensive and involves repeatedly resampling
  the data and fitting the compound model to histograms of each resampled
  dataset. The variance of the mean is estimated from the standard deviation of
  the returned mean parameters of every fit. 
\item[2.] \textbf{Weighted Gaussian and Linear Error Calculation}: \\
  For the Gaussian part of the compound mode, the error is calculated as 
  \[
    \varepsilon_\text{Gaussian} = \frac{\sigma_\text{Gaussian}}{\sqrt{N_\text{Gaussian} }}
  \] \vspace{3px}
  
  where $N_\text{Gaussian} $ is the number of stars within three standard deviations
  ($\sigma_\text{Gaussian} $) of the
  returned mean. 

  For the Linear part, the error is derived from the standard deviation of histogram bin residuals
  for bins outside the Gaussian region. The linear error is calculated as 

   \[
     \varepsilon_\text{Linear}
     = \frac{\sigma_\text{Linear}}{\sqrt{N_\text{Linear}}}
  \] \vspace{3px}
  
  where $N_\text{Linear} $ is the number of stars under the linear-part of the
  fit. 

  These errors are then weighted by the ratio of stars under each part of the
  model: 

  \[
  \varepsilon_\text{overall} = \sqrt{\verb|gaussian_weight| \cdot
  {\varepsilon_\text{Gaussian}}^2 + \verb|linear_weight| \cdot
{\varepsilon_\text{Linear}}^2}
  \] \vspace{3px}
  
  This error calculation is repeated for histogram bin numbers in the range of
  $\verb|int(|\sqrt{N_\text{total}}\verb|)|$ to
  $\verb|int(|\sqrt{N_\text{total}} + 10\verb|)|$. The minimum error from these
  calculations is chosen as the overall error. I have found this error almost
  nearly matches the error derived via bootstrapping for most cases. This runs
  much faster however. 
\end{itemize}

The entire OCF algorithm is repeated for initial \verb|catalogy| histogram bin
numbers ranging from eight to twenty. This is to ensure the optimal successful
bin parameters are found, as increasing the number of histogram bins affects
the generated compound models.

After completing all the compound fitting iterations, the algorithm selects the
parameters of the fit that correspond to the minimum error. These optimized
parameters provide an accurate and reliable mean magnitude for the starting
bin.

\subsubsection{Handling OCF Algorithm Failures} 

There is a possibility that the OCF algorithm might fail to find any successful
compound models, even after iterating through various amplitude values. The
primary reason for this failure is often the lack of sufficient stars in the
bin to fit an effective \verb|Gaussian1D() + Linear1D()| compound model. 

To address this issue, the algorithm includes a mechanism to extend the bin in
color-space and reattempt the fitting process: 

\begin{itemize}
  \item[1.] \textbf{Initial Failure Handling}: \\
    - If the OCF Algorithm fails for the starting bin, it extends the bin by 0.1
    in color-space of the CMD to include more stars. 

    - The algorithm then again extracts the stars in the new, larger bin and
    re-runs the OCF algorithm on the extended bin. 
  \item[2.] \textbf{Iterative Expansion}: \\
    - If the OCF algorithm fails again, the bin is extended another 0.1 in
    color-space. 

    - This process is repeated up to 7 times, each time extending the bin by 0.1,
    to include a progressively larger number of stars. 
  \item[3.] \textbf{Final Failure Condition}: \\
    - If the OCF algorithm to find a successful compound model after 7 extensions
    (a total extension of 0.7), it indicates a deeper issue.  

    - In this case, the failure is likely due to improper initial cutoffs
    ($x$-range / parallel line parameters). From my experience it is mainly
    from the parallel line parameters, after the 3x width increase, extending
    into the main-sequence part of the CMD, which prevents a single one-peak
    compound model from succeeding. 
\end{itemize} 

After completing the OCF algorithm for the starting bin, the algorithm proceeds
by generating a new bin starting at the end of the previous bin. This new bin
similarly begins having a width of \verb|x_range| / $n$, and OCF is again
implemented on the new bin. This iterative process continues, with each new bin
starting where the previous one ends, until the rightmost $x$ value of a bin
exceeds the end of the provided $x$-range. OCF is applied one last time to this
final bin and the algorithm concludes. 

After a successful OCF run across the RC cluster, the algorithm outputs a plot
like the one below for the example CMD: 

\begin{figure}[H]
  \centering
    \includegraphics[width = 9.8cm]{F115W-F212N-vsF115W_result-hist.png}
\end{figure}
\begin{figure}[H]
  \centering
    \includegraphics[width = 9.8cm]{F115W-F212N-vsF115W_result.png}
\end{figure}

The slope is finally calculated using a weighted
\verb|scipy.optimize.curve_fit()| linear fit which produces the following plot
in Figure 3: 

\begin{figure}[H]
  \centering
    \includegraphics[width = \textwidth]{F115W-F212N-vsF115W_slope.png}
    \caption{Calculated Slope using OCF algorithm for example CMD. Note that
    the Error Bars are scaled by 20x.}
\end{figure}

\subsection{Goodness-of-Fit Assesment for Entire Compound Model Fitting}

To ensure the accuracy and reliability of the derived RC (Red Clump) slope from
the multi-compound model fitting, we assess the goodness-of-fit by generating a 
synthetic RC cluster using the optimized compound models and performing
subsequent statistical tests. 

\begin{itemize}
  \item[1.] \textbf{Generating Synthetic RC Clusters}: \\
    - For each tile of the RC cluster, a synthetic RC cluster is generating
    using its corresponding optimized compound model. 
    
    - This synthetic cluster is designed to mimic the distribution of the
    actual RC tile. Example synthetic RC Tile clusters are shown in Figure 4. 
\item[2.] \textbf{Performing the Kolmogorov-Smirnov (KS) Test}: \\
    - The KS test is employed to statistically compare the entire collection of
    synthetic RC clusters against the actual RC cluster. 
    
    - This test evaluates the goodness-of-fit by calculating the distance
    (\verb|ks_statistic|)
    between the empirical
    cumulative distribution functions (CDFs) of the synthetic and actual RC
    clusters. 

    - A good fit, based on the KS test, thus has a \verb|ks_statistic| close to
    0. 
 \item[3.] \textbf{Plotting CDFs and Density Plots}: \\
   - To visually inspect the similarity between the synthetic and actual RC
   clusters, CDFs and density plots are generated. 

   - These plots help in identifying any significant deviations between the
   distributions of the synthetic and actual RC clusters. 

   - Example CDF and Density Plots are shown in Figure 5 and 6, respectively. 
\end{itemize}

When plotted using identical $x$-values, the synthetic cluster appears visually
distinct from the actual RC cluster. This difference arises because each tile
is generated from a single optimized compound model, resulting in a synthetic
cluster that only follows the characteristics (histogram) of that model rather than the
exact visual traits of the actual RC cluster. But for the purposes of
calculating the mean, the histograms closely resembling each other is what
matters. 

After the goodness-of-fit has been established, several statistical plots can
be generated: 

\subsubsection{Plotting Synthetic vs. Actual RC clusters on Example CMD} 

\begin{figure}[H]
  \centering
    \includegraphics[width = 11cm]{plot_hexbin.png}
    \caption{Hexbin plots comparing synthetic vs. actual RC tile data for tiles
    12 and 13 from an $n=18$ run.}
\end{figure}

\subsubsection{Plotting CDFs} 

\begin{figure}[H]
  \centering
    \includegraphics[width = 11cm]{plot_cdf.png}
    \caption{Overlaid synthetic \& actual CDFs for tiles 12 and 13 from an
    $n=18$ run.}
\end{figure}

\subsubsection{Plotting Densities} 

\begin{figure}[H]
  \centering
    \includegraphics[width = 11cm]{plot_kde.png}
    \caption{Overlaid density plots generated using a KDE for tiles 12 and 13
    from an $n=18$ run.}
\end{figure}


And finally, plotting the residuals which as a histogram, should resemble
a Gaussian.

\subsubsection{Plotting Residuals} 

\begin{figure}[H]
  \centering
    \includegraphics[width = \textwidth]{F115W-F212N-vsF115W_residual.png}
    \caption{Residuals from the calculated linear fit from the OCF algorithm
    from an $n=18$ run.}
\end{figure}

\subsection{Running the Algorithm} 

The algorithm requires an input parameter $n$, an integer telling the algorithm
how many tiles to segment the RC bar into. However, different values of $n$
could result in calculating different slopes. 

For this reason, the algorithm lets a user input an array of different $n$s
they wish to run the algorithm for. It then finds the optimal $n$ such that the
\textit{slope}-error is minimized and generates the above plots for the optimal
$n$ value. 

It also generates a plot like below, allowing one to compare the results of
different $n$ and any patterns that emerge across calculated slopes and errors
as $n$ increases. 

\begin{figure}[H]
  \centering
    \includegraphics[width = \textwidth]{F115W-F212N-vsF115W_[8, 9, 10, 11, 12, 13, 14, 15].png}
    \caption{Comparing OCF slope and error results across different $n$}
\end{figure}








\newpage
\section{Alternative OCF Algorithm -- Rectangles} 

For CMDs that look like this: 

\begin{figure}[H]
  \centering
    \includegraphics[width = \textwidth]{F212N-F323N-vsF212N_cutoff.png}
    \caption{NRCB1 F212N - F323N vs. F212 CMD}
\end{figure}

where it is potentially difficult to confidently generate the required parallel
line RC cutoff input parameters for the tiled-OCF algorithm detailed above,  there
is another option using an overall rectangular cutoff rather than two parallel
cutoffs.

This is fairly self-explanatory. Rather than generating tiled-bins that follow
the parallel lines across the slope of the RC, the algorithm generates
equal-height tiles that span the RC cluster. This method replaces the parallel
cutoff input parameters with an \verb|xlim| and \verb|ylim| parameter that
defines the overall rectangular RC cutoff within which the tiles are generated.
\\ 

The rest of the OCF algorithm is identical. \\

Note that this method is discouraged as the main method of calculating RC
slopes. I recommend using this method only to calculate a slope from which the
parallel cutoffs of the first method can better be determined. Usually though,
unsharp-masking is more than good enough.

In addition, more often than not, it is impossible to generate a perfect
rectangle RC cutoff using an \verb|xlim| and a \verb|ylim| that does not cross
into the main-sequence portion of the CMD (which would fail the OCF algorithm). 

\vfill

\href{https://github.com/devdeliw/research/tree/main/work}{\faGithub} --
\verb|red_clump_riemann.py|, \verb|red_clump_script.py|.







\end{document}

